\documentclass[11pt,a4paper]{article}
\usepackage[utf8]{inputenc}
\usepackage[T1]{fontenc}
\usepackage{amsmath,amssymb,amsthm}
\usepackage{geometry}
\usepackage{hyperref}
\usepackage{booktabs}
\usepackage{array}
\usepackage{listings}
\usepackage{xcolor}

\geometry{margin=2.8cm}

\theoremstyle{definition}
\newtheorem{definition}{Definition}[section]
\newtheorem{lemma}[definition]{Lemma}
\newtheorem{proposition}[definition]{Proposition}
\newtheorem{theorem}[definition]{Theorem}
\newtheorem{corollary}[definition]{Corollary}
\newtheorem{example}[definition]{Example}

\newcommand{\lt}{<}
\newcommand{\gt}{>}

\lstset{
  basicstyle=\ttfamily\small,
  columns=flexible,
  keepspaces=true,
  frame=single,
  backgroundcolor=\color{gray!10},
  rulecolor=\color{gray!50},
  breaklines=true,
  postbreak=\mbox{\textcolor{red}{$\hookrightarrow$}\space},
}

\newcommand{\Nat}{\mathbb{N}}
\newcommand{\val}{\operatorname{Val}}
\newcommand{\ones}[1]{1^{#1}}
\newcommand{\zeros}[1]{0^{#1}}
\newcommand{\List}{\operatorname{List}}
\newcommand{\conc}{{^\frown}}
\newcommand{\pre}{\text{pre}}
\newcommand{\foldl}{\operatorname{foldl}}
\newcommand{\symfun}{\operatorname{symFun}}
\newcommand{\compFun}{\operatorname{compFun}}
\newcommand{\config}{\operatorname{config}}
\newcommand{\countC}{\operatorname{countC}}
\newcommand{\auxA}{\mathcal{A}}
\newcommand{\auxB}{\mathcal{B}}
\newcommand{\auxRules}{\mathcal{X}}
\newcommand{\dynRules}{\mathcal{D}}
\newcommand{\antihydraSRS}{\mathcal{R}_{A}}
\newcommand{\macroStep}{\operatorname{Step}}
\newcommand{\nextMathState}{\operatorname{nextMathState}}
\newcommand{\mathHalts}{\operatorname{mathHalts}}
\newcommand{\some}{\operatorname{some}}
\newcommand{\none}{\operatorname{none}}
\newcommand{\valOrbit}{\operatorname{valOrbit}}
\newcommand{\counterZ}{\operatorname{counterZ}}
\newcommand{\grams}{\operatorname{grams}}
\newcommand{\lpad}{\texttt{\^{}}}
\newcommand{\rpad}{\texttt{\$}}
\newcommand{\macroIter}{\operatorname{macroIter}}
\newcommand{\Option}{\operatorname{Option}}
\newcommand{\bind}{\mathbin{>\!\!\!>\!=}}
\newcommand{\futureDip}{\operatorname{futureDip}}
\newcommand{\orbitCounter}{\operatorname{orbitCounter}}
\newcommand{\IsHalting}{\operatorname{IsHalting}}
\newcommand{\RegularInvariant}{\operatorname{RegularInvariant}}
\newcommand{\ZZ}{\mathbb{Z}}
\newcommand{\getD}{\operatorname{getD}}

\title{A mixed-base encoding and FAR-style invariant search for Antihydra}
\author{}
\date{}

\begin{document}
\maketitle

\begin{abstract}
  Antihydra is the BB(6) candidate
  \texttt{1RB1RA\_0LC1LE\_1LD1LC\_1LA0LB\_1LF1RE\_---0RA}.
  We give a self-contained account of its integer dynamics and of the
  mixed binary--ternary rewriting system used in the FAR/CPS search.
  The rewriting system is derived from the macro-level analysis of the
  underlying Turing machine; its transport rules preserve value and
  counter, and every rewriting strategy simulates the integer orbit
  step-for-step.
\end{abstract}

\section{From the Turing machine to integer dynamics}
\label{sec:dynamics}

\subsection{The Antihydra Turing machine and its macro rules}
\label{subsec:macro}

Antihydra is the 6-state Turing machine with the following transition
 table.  The tape alphabet is $\{0,1\}$ and the states are
$\{A,B,C,D,E,F\}$; the initial state is $A$ and the initial tape is
blank.

\begin{center}
\begin{tabular}{c|cc}
  \toprule
  state & read $0$ & read $1$ \\
  \midrule
  $A$ & $1RB$ & $1RA$ \\
  $B$ & $0LC$ & $1LE$ \\
  $C$ & $1LD$ & $1LC$ \\
  $D$ & $1LA$ & $0LB$ \\
  $E$ & $1LF$ & $1RE$ \\
  $F$ & halt  & $0RA$ \\
  \bottomrule
\end{tabular}
\end{center}

For the analysis it is convenient to work with a macro-level
configuration.  Following the Lean formalisation in
\texttt{bbchallenge/Antihydra/Antihydra.lean}, for
$b,m,n,p\in\Nat$ define
\begin{equation}
  \label{eq:macro-config}
  P(b,m,n,p)
  =
  \bigl(\text{state }E,\;\text{head }0,\;
  \text{left}=\ones{m}\,0\,\ones{b},\;
  \text{right}=\ones{n}\,\zeros{p}\bigr).
\end{equation}
The block $\ones{m}$ is the ``value block'', the block $\ones{b}$ is
a unary counter, and the right block $\ones{n}$ is used during the
macro step.

The machine-checked macro theorems proved in \texttt{Antihydra.lean}
are the following.  (We write $C\xrightarrow{t}D$ to mean that the
configuration $C$ reaches $D$ after exactly $t$ steps.)

\begin{lemma}[Macro loop step]
  \label{lem:p-step}
  For all $b,m,n,p\in\Nat$,
  \[
    P\bigl(b,m+4,n,p+2\bigr)
    \xrightarrow{\,2n+12\,}
    P\bigl(b,m+2,n+3,p+1\bigr).
  \]
\end{lemma}

\begin{lemma}[Macro loop, $k$ iterations]
  \label{lem:p-multistep}
  For all $b,m,n,p,k\in\Nat$,
  \[
    P\bigl(b,m+2+2k,n,p+1+k\bigr)
    \xrightarrow{\,k(2n+3k+9)\,}
    P\bigl(b,m+2,n+3k,p+1\bigr).
  \]
\end{lemma}

\begin{lemma}[Even endgame]
  \label{lem:even-endgame}
  For all $b,N,p\in\Nat$, the configuration $P(b,0,N,p+2)$ halts in
  exactly $2$ steps.
\end{lemma}

\begin{lemma}[Odd endgame, counter positive]
  \label{lem:odd-endgame}
  For all $b',N,p\in\Nat$,
  \[
    P\bigl(b'+1,3,N,p+2\bigr)
    \xrightarrow{\,3N+20\,}
    P\bigl(b',N+6,0,p\bigr).
  \]
\end{lemma}

\begin{lemma}[Even endgame to next loop]
  \label{lem:even-to-loop}
  For all $b,N,p\in\Nat$,
  \[
    P\bigl(b,2,N,p+2\bigr)
    \xrightarrow{\,3N+2b+26\,}
    P\bigl(b+2,N+5,0,p\bigr).
  \]
\end{lemma}

\begin{lemma}[Odd halt endgame]
  \label{lem:odd-halt}
  For all $N,p\in\Nat$, the configuration $P(0,3,N,p+2)$ halts in
  exactly $2N+12$ steps.
\end{lemma}

These six facts are enough to reduce the behaviour of the machine,
once it has entered the macro regime, to a discrete dynamical system on
a pair of nonnegative integers.

\subsection{The integer recurrence}
\label{subsec:recurrence}

Projecting the macro transitions onto the appropriate integer value
coordinate gives the following familiar Collatz-like map.  Write the
integer value as $a\in\Nat$ and the counter as $b\in\Nat$.  The
macro analysis implies that, as long as $a\ge 2$, the next pair is
\begin{equation}
  \label{eq:antihydra-map}
  (a,b)\longmapsto
  \begin{cases}
    \bigl(\lfloor 3a/2\rfloor,\;b+2\bigr), & a\text{ even},\\[4pt]
    \bigl(\lfloor 3a/2\rfloor,\;b-1\bigr), & a\text{ odd and }b>0,\\[4pt]
    \text{halt}, & a\text{ odd and }b=0.
  \end{cases}
\end{equation}
Equivalently, writing $a=2n$ when $a$ is even and $a=2n+1$ when $a$ is
odd,
\begin{equation}
  \label{eq:antihydra-map-explicit}
  (2n,b)\mapsto(3n,b+2),
  \qquad
  (2n+1,b)\mapsto
  \begin{cases}
    (3n+1,b-1), & b>0,\\
    \text{halt}, & b=0.
  \end{cases}
\end{equation}
The machine halts precisely when an odd value is encountered while the
counter is empty.  Antihydra is believed to run forever, because the
walk on the counter has positive drift ($+2$ with probability $1/2$,
$-1$ with probability $1/2$).

The initial macro configuration reached after the $58$-step warm-up is
$P(2,8,0,0)$.  In the mixed-base encoding introduced below this
configuration is written $\langle fff\rangle$ and has value $a_0=8$ and
counter $b_0=0$.

\paragraph{Remark on the value coordinate.}
In \texttt{Antihydra.lean} the abstract \texttt{MathState} tracks the
length $m$ of the left block $\ones{m}$ in
\eqref{eq:macro-config}; its even branch reads $2n\mapsto 3n+2$.
The mixed-base encoding below uses a different Turing-machine invariant,
namely the value $\val(\lt w\gt)$ of the digit string, which satisfies
the cleaner Mahler form \eqref{eq:antihydra-map}.  Both invariants give
the same halting condition, and the Python simulation in
\texttt{antihydra\_far.py} validates \eqref{eq:antihydra-map} against the
actual Turing-machine orbit for hundreds of dynamic steps.

\section{Mixed-base rewriting rules}
\label{sec:rewriting}

\subsection{Alphabet and value function}
\label{subsec:alphabet}

The mixed binary--ternary encoding uses the alphabet
\[
  \Sigma=\{\;\lt\, ,\; \gt\, ,\; f,\; t,\; 0,\; 1,\; 2,\; c\;\}.
\]
The symbols $f$ and $t$ are binary digits ($f$ for $2x$, $t$ for
$2x+1$); the symbols $0,1,2$ are ternary digits ($3x$, $3x+1$,
$3x+2$).  The symbol $\lt$ is the left boundary, $\gt$ is the right
boundary of the value word, and $c$ is a unary counter symbol.

A \emph{configuration} is a string of the form
\begin{equation}
  \label{eq:config}
  \lt w\gt c^{b},
  \qquad w\in\{f,t,0,1,2\}^{*},\; b\in\Nat.
\end{equation}
The value of the configuration is computed by reading the value word
$\lt w\gt$ from left to right:
\begin{equation}
  \label{eq:val}
  \val(\lt w\gt c^{b})
  =
  \val_{\lt}(w),
\end{equation}
where the partial function $\val_{x}(u)$ is defined recursively by
\begin{alignat}{2}
  \val_{x}(\varepsilon) &= x,\nonumber\\
  \val_{x}(f\,u) &= \val_{2x}(u),\nonumber\\
  \val_{x}(t\,u) &= \val_{2x+1}(u),\nonumber\\
  \val_{x}(d\,u) &= \val_{3x+d}(u)\qquad (d\in\{0,1,2\}),\nonumber\\
  \val_{x}(\gt\,u) &= x.\label{eq:val-rec}
\end{alignat}
Thus $\lt$ contributes the initial value $1$, each subsequent digit
applies the corresponding affine function, and $\gt$ returns the
current value.  The counter symbols $c$ do not affect the value.

\begin{example}
  The initial configuration is $\lt fff\gt$.  Its value is
  \[
    \val(\lt fff\gt)
    = 1\xrightarrow{f}2\xrightarrow{f}4\xrightarrow{f}8.
  \]
  More generally, a binary word $w$ ending in $f$ or $t$ represents the
  integer whose binary expansion is the sequence of bits encoded by
  $w$.
\end{example}

\subsection{The rewriting system}
\label{subsec:rules}

The rewriting relation is generated by the following rules.  The
rules are divided into three groups: dynamic rules, which implement the
integer map and change the counter; transport rules, which are
value-preserving and move digits; and boundary rules, which turn a
ternary digit at the left end back into binary digits.

\paragraph{Dynamic rules.}
\begin{equation}
  \label{eq:dyn}
  \begin{aligned}
    f\gt &\;\longrightarrow\; 0\gt cc,
    &&\text{(even value: }2n\mapsto 3n,\; b\mapsto b+2\text{)},\\
    t\gt c &\;\longrightarrow\; 1\gt,
    &&\text{(odd value: }2n+1\mapsto 3n+1,\; b\mapsto b-1\text{)}.
  \end{aligned}
\end{equation}
The second rule is applicable only when a counter symbol $c$ is
present; if the configuration ends in $t\gt$ with no $c$ then no rule
applies and the configuration is halting.

\paragraph{Transport rules.}
\begin{equation}
  \label{eq:transport}
  \begin{array}{lll}
    f0\to 0f,\qquad & f1\to 0t,\qquad & f2\to 1f,\\
    t0\to 1t,\qquad & t1\to 2f,\qquad & t2\to 2t.
  \end{array}
\end{equation}
Each of these rules preserves the represented value: for example,
$f0(x)=2\cdot(3x)=6x=0f(x)$ and $t1(x)=3\cdot(2x+1)+1=6x+4=2f(x)$.

\paragraph{Boundary rules.}
\begin{equation}
  \label{eq:boundary}
  \lt 0\to\lt t,\qquad
  \lt 1\to\lt ff,\qquad
  \lt 2\to\lt ft.
\end{equation}
These implement the equation $3=2\cdot1+1$, $4=2\cdot(2\cdot1)$, and
$5=2\cdot(2\cdot1)+1$ at the left boundary.

The full set of rules is therefore
\begin{equation}
  \label{eq:all-rules}
  \begin{aligned}
  \mathcal{R}
  ={}&\{f\gt\to0\gt cc,\; t\gt c\to1\gt\}\\
  \cup{}&\{f0\to0f,\; f1\to0t,\; f2\to1f,\; t0\to1t,\; t1\to2f,\; t2\to2t\}\\
  \cup{}&\{\lt0\to\lt t,\; \lt1\to\lt ff,\; \lt2\to\lt ft\}.
  \end{aligned}
\end{equation}

\begin{definition}[Initial and halting configurations]
  The initial configuration is $\lt fff\gt$, which has value $a_0=8$
  and counter $b_0=0$.  A configuration is \emph{halting} if it ends in
  $t\gt$ and contains no counter symbol, i.e. it has the form
  $u\,t\gt$ with $u\in\Sigma^{*}$ and no $c$ occurs.
\end{definition}

\subsection{Simulation}
\label{subsec:simulation}

The rewriting system is designed so that the rightmost-rule strategy
simulates the integer recurrence \eqref{eq:antihydra-map}.  More
precisely, define a \emph{dynamic firing} to be an application of one
of the two dynamic rules.  Then:

\begin{itemize}
\item Every transport or boundary rule preserves both the value
  $\val(\cdot)$ and the number of counter symbols.
\item A dynamic firing $f\gt\to0\gt cc$ applies exactly when the
  current value $a$ is even, and it changes the value to
  $\lfloor3a/2\rfloor$ while adding two counter symbols.
\item A dynamic firing $t\gt c\to1\gt$ applies exactly when the
  current value $a$ is odd and the counter is positive, and it changes
  the value to $\lfloor3a/2\rfloor$ while removing one counter symbol.
\item If the value is odd and the counter is empty, the configuration
  ends in $t\gt$ and no rule applies: this is the halting condition.
\end{itemize}

A useful structural fact, inherited from the YAH framework, is that
the counter available at the $j$-th odd step is independent of the
rewriting strategy: transport and boundary rules preserve the value
and the counter, and dynamic steps depend only on the parity of the
current value.  This is the analogue for Antihydra of the
strategy-independence argument in the proof of Theorem~3.17, Claim~1
of Yolcu--Aaronson--Heule~(\cite[pp.~20--21]{YAH21}): there the
auxiliary rules $X=A\cup B$ preserve $\val(\cdot)$ while the dynamic
rules $DT$ apply the Collatz map $T$ to the value.  Hence the
nondeterministic closure of the rewriting system decides exactly the
same halting question as the deterministic integer orbit.

\begin{example}
  The first few dynamic steps from the initial configuration are:
  \[
    \lt fff\gt
    \xrightarrow{f\gt\to0\gt cc}
    \lt ff0\gt cc
    \xrightarrow{\text{transport}}
    \cdots
    \xrightarrow{f\gt\to0\gt cc}
    \lt 0\dots0\gt c^{b}.
  \]
  At each dynamic firing the represented value follows
  $8,12,18,27,\dots$, and the counter follows $0,2,4,6,\dots$,
  exactly as prescribed by \eqref{eq:antihydra-map}.
\end{example}

\section{Formal verification in Lean}
\label{sec:lean}

\begin{sloppypar}
The claims of the preceding sections have been formalised in Lean~4 in
the directory \texttt{lean-code/SRS}.  The relevant files are
\texttt{Basic.lean} (string-rewriting systems),
\texttt{MixedBaseRepresentation.lean} (the YAH mixed-base value
function), \texttt{AntihydraSRSaxiom.lean} (a BusyLean-free fragment of
the Antihydra macro model), \texttt{AntihydraMahler.lean} (the equivalent
Mahler-form value coordinate \eqref{eq:antihydra-map} and its agreement
with the macro model), \texttt{AntihydraSRS.lean} (the SRS itself
and its local correctness proof), and
\texttt{AntihydraSRSSimulation.lean} (the global faithful simulation
between arbitrary SRS derivations and the deterministic macro orbit), and
\texttt{AntihydraSRSForward.lean} (a constructive forward strategy that
realises the orbit, giving the exact converse of the faithful-simulation
theorem), and
\texttt{AntihydraSRSObstruction.lean} (a pumping obstruction showing that
unbounded future dips of the counter rule out any regular non-halting
invariant over unary-counter encodings).
This section records the precise definitions and the main lemmas with
proofs in mathematical prose.
\end{sloppypar}

\subsection{String rewriting systems}

Let $\alpha$ be a type of symbols.  Following \texttt{Basic.lean}, a
\emph{string rewriting system} (SRS) over $\alpha$ is a binary relation
$R$ on $\List\alpha$; we write $R\,\ell\,r$ or $\ell\to r\in R$ when
$(\ell,r)$ is a rule.  The \emph{rewrite relation} $\to_R$ is defined
by
\begin{equation}
  \label{eq:rewrite-step}
  u\to_R v
  \;\iff\;
  \exists\,s,t,\ell,r\;\bigl(\ell\to r\in R\;\wedge\;u=s\ell t\;\wedge\;v=srt\bigr).
\end{equation}
In words: a rule may be applied inside any surrounding context.  The
following two closure properties are immediate from the definition.

\begin{lemma}[Rules are rewrite steps]
  If $\ell\to r\in R$ then $\ell\to_R r$.
\end{lemma}

\begin{lemma}[Context closure]
  If $u\to_R v$ then $sut\to_R svt$ for all strings $s,t$.
\end{lemma}

For systems $R,S$ we write $R\cup S$ for the union of their rule sets.

\subsection{Mixed-base representations}
\label{subsec:mixed-base}

A \emph{mixed digit} is a pair $(n)_b$ with digit $n\in\Nat$ and base
$b\in\Nat$; a string of such digits is a mixed-base representation.  In
the function view of \texttt{MixedBaseRepresentation.lean} the symbol
$(n)_b$ denotes the affine map
\[
  \beta_b^n(x)=b\cdot x+n,
\]
and a string $N=(n_1)_{b_1}\cdots(n_k)_{b_k}$ denotes the composite
\begin{equation}
  \label{eq:Gamma}
  \Gamma_N(x)
  =
  (\beta_{b_k}^{n_k}\circ\cdots\circ\beta_{b_1}^{n_1})(x),
\end{equation}
where the leftmost symbol is applied innermost.  The \emph{value} of the
string is the natural number
\begin{equation}
  \label{eq:val-mixed}
  \val(N)=\sum_{i=1}^{k} n_i\prod_{j=i+1}^{k} b_j,
\end{equation}
with the convention that the empty product is $1$.

\begin{lemma}[\texttt{gamma\_eq}]
  For every mixed-base string $N$ and every $x\in\Nat$,
  \[
    \Gamma_N(x)=\bigl(\prod_{j=1}^{k} b_j\bigr)\cdot x+\val(N).
  \]
\end{lemma}
\begin{proof}
  By induction on the length of $N$.  The empty string gives
  $\Gamma_{\emptyset}(x)=x=1\cdot x+0$.  For $N=(n)_b::N'$, the
  induction hypothesis gives
  $\Gamma_{N'}(y)=(\prod_{j=2}^{k} b_j)\cdot y+\val(N')$.  Taking
  $y=\beta_b^n(x)=bx+n$ and using the definition of $\val$ yields the
  claim after collecting the $x$-term and the constant term.
\end{proof}

\begin{lemma}[\texttt{gamma\_headBase\_zero}]
  If the leading base is $b_1=0$, then $\Gamma_N(x)=\val(N)$ for every
  $x$.
\end{lemma}
\begin{proof}
  In \eqref{eq:Gamma} the innermost map is $\beta_0^{n_1}(x)=n_1$, so
  the whole composite is constant; by the previous lemma the constant is
  $\val(N)$.
\end{proof}

The fundamental local operation is \emph{base swapping}: two adjacent
symbols can have their bases exchanged without changing the composite,
at the cost of updating the digits by division with remainder.

\begin{lemma}[\texttt{baseSwap}]
  \label{lem:base-swap}
  Let $b,c,n,m\in\Nat$ with $n<b$ and $m<c$.  Define
  \[
    n'=\frac{bm+n}{c},
    \qquad
    m'=(bm+n)\bmod c.
  \]
  Then $n'<b$, $m'<c$, and for every $x$,
  \[
    \beta_b^n\bigl(\beta_c^m(x)\bigr)
    =
    \beta_c^{m'}\bigl(\beta_b^{n'}(x)\bigr).
  \]
\end{lemma}
\begin{proof}
  Both sides equal $bcx+bm+n$.  On the right we use the division
  identity $bm+n=cn'+m'$.  The bound $m'<c$ follows from the definition
  of the remainder.  For $n'<b$, note that $m<c$ implies
  $b(m+1)\le bc$, so $bm+n<bm+b\le bc$; hence $(bm+n)/c<b$.
\end{proof}

\subsection{The Antihydra alphabet and rules}
\label{subsec:lean-alphabet}

The Antihydra SRS uses the eight-symbol alphabet
\[
  \Sigma_{A}=\{f,t,0,1,2,\lt,\gt,c\},
\]
where in the Lean code the ternary digits are named $d0,d1,d2$ and the
boundary markers are $lhd,rhd$.  The symbol maps are
\begin{equation}
  \label{eq:symfun}
  \begin{array}{c|c}
    \text{symbol }s & \symfun_s(x) \\ \hline
    f & 2x \\
    t & 2x+1 \\
    0 & 3x \\
    1 & 3x+1 \\
    2 & 3x+2 \\
    \lt & 1 \\
    \gt & x \\
    c & x
  \end{array}
\end{equation}
Thus $\lt$ contributes the initial value $1$; $\gt$ returns the current
value; and the counter symbol $c$ is the identity.  For a word
$w\in\List{\Sigma_A}$ and a starting value $x$, the composite map is the
left fold
\[
  \compFun_w(x)=w.\foldl\bigl((acc,s)\mapsto\symfun_s(acc)\bigr)(x).
\]

\begin{lemma}[\texttt{compFun\_cons} and \texttt{compFun\_append}]
  \label{lem:compfun-props}
  For all symbols $s$, words $u,v$, and values $x$,
  \[
    \compFun_{s::w}(x)=\compFun_w\bigl(\symfun_s(x)\bigr)
  \]
  and
  \[
    \compFun_{u\conc v}(x)=\compFun_v\bigl(\compFun_u(x)\bigr).
  \]
\end{lemma}
\begin{proof}
  Both are immediate from the definition of a left fold and the
  associativity of list append.
\end{proof}

A \emph{configuration} is a string of the form
\begin{equation}
  \label{eq:lean-config}
  \config w b
  =
  \lt\;w\;\gt\;c^{b},
\end{equation}
where $w\in\{f,t,0,1,2\}^{*}$ is the value word and $b\in\Nat$ is the
counter.  The value of the configuration is independent of the counter.

\begin{lemma}[\texttt{cfgVal}]
  \label{lem:cfg-val}
  For every value word $w$, counter $b$, and value $x\in\Nat$,
  \[
    \val(\config w b)=\compFun_{\lt w\gt}(x).
  \]
\end{lemma}
\begin{proof}
  Since $\symfun_c(x)=x$, the block $c^{b}$ is transparent:
  \texttt{compFun\_replicate\_c} proves $\compFun_{c^{b}}(x)=x$ by
  induction on $b$.  Lemma~\ref{lem:compfun-props} therefore gives
  \[
    \compFun_{\lt w\gt c^{b}}(x)
    =\compFun_{c^{b}}\bigl(\compFun_{\lt w\gt}(x)\bigr)
    =\compFun_{\lt w\gt}(x).\qedhere
  \]
\end{proof}

The counter is measured by the number of $c$ symbols:
\[
  \countC(w)=\text{the number of occurrences of }c\text{ in }w.
\]
It is additive over concatenation: $\countC(uv)=\countC(u)+\countC(v)$.

The rewriting system is the union of three groups of rules.

\paragraph{Dynamic rules.}
\begin{equation}
  \label{eq:lean-dyn}
  \dynRules:\quad
  f\gt\to 0\gt cc,
  \qquad
  t\gt c\to 1\gt.
\end{equation}

\paragraph{Transport rules.}
\begin{equation}
  \label{eq:lean-transport}
  \auxA:\quad
  \begin{array}{lll}
    f0\to 0f, & f1\to 0t, & f2\to 1f, \\
    t0\to 1t, & t1\to 2f, & t2\to 2t.
  \end{array}
\end{equation}

\paragraph{Boundary rules.}
\begin{equation}
  \label{eq:lean-boundary}
  \auxB:\quad
  \lt 0\to\lt t,
  \qquad
  \lt 1\to\lt ff,
  \qquad
  \lt 2\to\lt ft.
\end{equation}

The auxiliary subsystem is $\auxRules=\auxA\cup\auxB$, and the full
Antihydra SRS is
\[
  \antihydraSRS=\dynRules\cup\auxRules.
\]

\subsection{The macro model}
\label{subsec:lean-macro}

To state what the SRS represents, we first recall the integer recurrence
from Section~\ref{subsec:recurrence} in the mixed-base value coordinate.
Define the (partial) \emph{macro step}
$\macroStep:\Nat\times\Nat\rightharpoonup\Nat\times\Nat$ by
\begin{equation}
  \label{eq:lean-macro-step}
  \macroStep(a,b)=
  \begin{cases}
    \bigl(\lfloor 3a/2\rfloor,\;b+2\bigr), & a\text{ even},\\
    \text{undefined}, & a\text{ odd and }b=0,\\
    \bigl(\lfloor 3a/2\rfloor,\;b-1\bigr), & a\text{ odd and }b>0.
  \end{cases}
\end{equation}
Equivalently, writing $a=2x$ or $a=2x+1$,
\begin{equation}
  \label{eq:lean-macro-explicit}
  \macroStep(2x,b)=(3x,b+2),
  \qquad
  \macroStep(2x+1,0)=\text{undefined},
  \qquad
  \macroStep(2x+1,b+1)=(3x+1,b).
\end{equation}

In the file \texttt{AntihydraSRSaxiom.lean} the same machine is
presented in a different, but equivalent, coordinate: the state is a
pair $(a,b)$ where $a$ is the length $m$ of the left block $\ones{m}$ in
the macro configuration $P(b,m,n,p)$ of \eqref{eq:macro-config}, and the
total map is
\[
  A(a,b)=\bigl(3\cdot(a/2)+2,\;b+2\bigr)\quad\text{if $a$ even},
  \qquad
  A(a,b)=\bigl(3\cdot(a/2)+3,\;b-1\bigr)\quad\text{if $a$ odd}.
\]
The partial step $\nextMathState$ returns $\none$ exactly when $a$ is
odd and $b=0$.  The reason for the $+2$/$+3$ offsets is that this model
tracks the length $m$ directly, whereas the mixed-base model tracks the
value $\val(\lt w\gt)$.  Both models agree on the halting condition, and
the bridge is supplied by the BusyLean proof in
\texttt{AntihydraMachine.lean} (see
Section~\ref{subsec:tm-bridge}).  Our SRS correctness theorem is stated
in the cleaner mixed-base coordinate.

\subsection{Correctness: the SRS represents the macro step}
\label{subsec:lean-correctness}

The correctness proof has three ingredients: the auxiliary rules
preserve value and counter, the dynamic rules implement the two branches
of $\macroStep$, and the halting condition matches.

\begin{lemma}[\texttt{auxRules\_preserve\_value}]
  \label{lem:aux-value}
  For every rule $\ell\to r\in\auxRules$ and every $x\in\Nat$,
  $\compFun_\ell(x)=\compFun_r(x)$.
\end{lemma}
\begin{proof}
  There are nine rules.  Recall that $\compFun$ is a left fold, so a
  string $s_1s_2\dots s_k$ means: start with $x$, apply $s_1$, then
  $s_2$, and so on, finishing with $s_k$; the final $\gt$ returns the
  current accumulator.

  \paragraph{Transport rules.}  The six transport rules are the special
  cases of Lemma~\ref{lem:base-swap} for bases $2$ and $3$.  Writing the
  left-fold computation in each case gives
  \[
    \begin{aligned}
      f0(x)&=3\cdot(2x)=6x,            & 0f(x)&=2\cdot(3x)=6x,\\
      f1(x)&=3\cdot(2x)+1=6x+1,        & 0t(x)&=2\cdot(3x)+1=6x+1,\\
      f2(x)&=3\cdot(2x)+2=6x+2,        & 1f(x)&=2\cdot(3x+1)=6x+2,\\
      t0(x)&=3\cdot(2x+1)=6x+3,        & 1t(x)&=2\cdot(3x+1)+1=6x+3,\\
      t1(x)&=3\cdot(2x+1)+1=6x+4,      & 2f(x)&=2\cdot(3x+2)=6x+4,\\
      t2(x)&=3\cdot(2x+1)+2=6x+5,      & 2t(x)&=2\cdot(3x+2)+1=6x+5.
    \end{aligned}
  \]
  In every pair the two composites agree.

  \paragraph{Boundary rules.}  Here the leftmost symbol $\lt$ replaces
  the input by $1$, after which the remaining digits are applied in order
  and $\gt$ returns the result.  Thus
  \[
    \begin{aligned}
      \compFun_{[\lt,0]}(x)&=0(1)=3,        &
      \compFun_{[\lt,t]}(x)&=t(1)=2\cdot1+1=3,\\
      \compFun_{[\lt,1]}(x)&=1(1)=4,        &
      \compFun_{[\lt,f,f]}(x)&=f(f(1))=2\cdot2=4,\\
      \compFun_{[\lt,2]}(x)&=2(1)=5,        &
      \compFun_{[\lt,f,t]}(x)&=t(f(1))=2\cdot2+1=5.
    \end{aligned}
  \]
  All three boundary rules therefore preserve value.
\end{proof}

\begin{lemma}[\texttt{auxRules\_preserve\_count}]
  \label{lem:aux-count}
  Every rule $\ell\to r\in\auxRules$ satisfies $\countC(\ell)=\countC(r)$.
\end{lemma}
\begin{proof}
  None of the nine auxiliary rules mentions the counter symbol $c$.
\end{proof}

Because rewriting is applied inside a context, preservation extends from
rules to arbitrary rewrite steps.

\begin{theorem}[\texttt{auxRules\_rewriteStep\_preserve\_value} and
  \texttt{\_preserve\_count}]
  \label{thm:aux-step}
  If $u\to_{\auxRules}v$, then $\compFun_u(x)=\compFun_v(x)$ for all
  $x$ and $\countC(u)=\countC(v)$.
\end{theorem}
\begin{proof}
  Write $u=s\ell t$ and $v=srt$ with $\ell\to r\in\auxRules$.  Put
  $y=\compFun_s(x)$.  By Lemma~\ref{lem:compfun-props},
  \[
    \compFun_u(x)=\compFun_t\bigl(\compFun_\ell(y)\bigr)
    \:\text{and}\:
    \compFun_v(x)=\compFun_t\bigl(\compFun_r(y)\bigr),
  \]
  which are equal by Lemma~\ref{lem:aux-value}.  The counter equality
  follows from $\countC(uv)=\countC(u)+\countC(v)$ and
  Lemma~\ref{lem:aux-count}.
\end{proof}

\begin{lemma}[\texttt{dynEven}]
  \label{lem:dyn-even}
  For every prefix word $\pre$ and every $x\in\Nat$,
  \begin{align*}
    \compFun_{\pre\conc f\gt}(x)&=2\cdot\compFun_\pre(x),\\
    \compFun_{\pre\conc 0\gt cc}(x)&=3\cdot\compFun_\pre(x),\\
    \countC(\pre\conc 0\gt cc)&=\countC(\pre\conc f\gt)+2.
  \end{align*}
\end{lemma}
\begin{proof}
  Direct computation.  We have
  $\compFun_{\pre\conc f\gt}(x)=\gt\bigl(\symfun_f(\compFun_\pre(x))\bigr)
  =2\cdot\compFun_\pre(x)$, and similarly
  $\compFun_{\pre\conc 0\gt cc}(x)=\gt\bigl(\symfun_0(\compFun_\pre(x))\bigr)
  =3\cdot\compFun_\pre(x)$ because the trailing $cc$ is transparent.
  The counter changes by exactly the two new $c$ symbols.
\end{proof}

\begin{lemma}[\texttt{dynOdd}]
  \label{lem:dyn-odd}
  For every prefix word $\pre$ and every $x\in\Nat$,
  \begin{align*}
    \compFun_{\pre\conc t\gt c}(x)&=2\cdot\compFun_\pre(x)+1,\\
    \compFun_{\pre\conc 1\gt}(x)&=3\cdot\compFun_\pre(x)+1,\\
    \countC(\pre\conc t\gt c)&=\countC(\pre\conc 1\gt)+1.
  \end{align*}
\end{lemma}
\begin{proof}
  Same direct computation, using $\symfun_t(y)=2y+1$ and
  $\symfun_1(y)=3y+1$.
\end{proof}

Together, Lemmas \ref{lem:dyn-even} and \ref{lem:dyn-odd} say that a
dynamic firing on a configuration changes the value word from a binary
least significant digit $f$ or $t$ to a ternary least significant digit
$0$ or $1$, while applying exactly the map $a\mapsto\lfloor3a/2\rfloor$
and updating the counter by $+2$ or $-1$.

\begin{theorem}[\texttt{antihydraSRS\_represents\_macro}]
  \label{thm:srs-represents-macro}
  The Antihydra SRS represents the macro step $\macroStep$ in the
  following sense.
  \begin{enumerate}
  \item Every auxiliary rewrite preserves value and counter:
    if $u\to_{\auxRules}v$ then
    $\compFun_u(x)=\compFun_v(x)$ for all $x$ and
    $\countC(u)=\countC(v)$.
  \item The even dynamic rule realises
    $(2x,b)\mapsto(3x,b+2)$: for every prefix $\pre$,
    \begin{align*}
      \compFun_{\pre\conc f\gt}(x)&=2\cdot\compFun_\pre(x),\\
      \compFun_{\pre\conc 0\gt cc}(x)&=3\cdot\compFun_\pre(x),
    \end{align*}
    and the counter grows by $2$.
  \item The odd dynamic rule realises
    $(2x+1,b+1)\mapsto(3x+1,b)$: for every prefix $\pre$,
    \begin{align*}
      \compFun_{\pre\conc t\gt c}(x)&=2\cdot\compFun_\pre(x)+1,\\
      \compFun_{\pre\conc 1\gt}(x)&=3\cdot\compFun_\pre(x)+1,
    \end{align*}
    and the counter drops by $1$.
  \end{enumerate}
\end{theorem}
\begin{proof}
  Part~(1) is Theorem~\ref{thm:aux-step}; parts~(2) and~(3) are
  Lemmas \ref{lem:dyn-even} and \ref{lem:dyn-odd}.
\end{proof}

\subsection{Faithful simulation}
\label{subsec:lean-simulation}

Theorem~\ref{thm:srs-represents-macro} shows that each rule of the SRS
does the right thing \emph{when it fires}.  It does not yet say that an
arbitrary SRS derivation must follow the deterministic macro orbit: a
derivation might, in principle, fire dynamic rules at unexpected
positions or in an order inconsistent with the value.  The file
\texttt{AntihydraSRSSimulation.lean} closes this gap by proving that
\emph{every} derivation from a configuration is faithful to the orbit.

A \emph{digit string} is a word
$w\in\{f,t,0,1,2\}^{*}$ containing no markers or counter symbol.  A
\emph{configuration} is therefore $\config w b$ with $w$ a digit string.
The first structural fact is that a dynamic redex is forced to the
unique boundary marker $\gt$.

\begin{lemma}[\texttt{dynStep\_config\_dest}]
\label{lem:dyn-align}
Let $w$ be a digit string and $b\in\Nat$.  If
$\config w b\to_{\dynRules}v$, then either
\begin{itemize}
\item $w=w'\conc f$ and $v=\config{(w'\conc0)}{(b+2)}$, or
\item $w=w'\conc t$, $b=b'+1$, and $v=\config{(w'\conc1)}{b'}$.
\end{itemize}
\end{lemma}
\begin{proof}[Proof sketch]
A dynamic rule contains the marker $\gt$ in its left-hand side, while a
configuration contains exactly one $\gt$ (because $w$ is a digit
string).  By counting occurrences of $\gt$ and using injectivity of the
split around that symbol, the redex must be positioned so that the rule's
$\gt$ coincides with the configuration's unique $\gt$.  Hence the rule
sees the last symbol of $w$; if it is $f$ the even rule fires, if it is
$t$ the odd rule fires and consumes one counter symbol.
\end{proof}

Consequently the rule that fires is determined by the parity of the
value $\val(\config w b)$, not by the rewriting strategy.

\begin{lemma}[\texttt{dynStep\_config}]
\label{lem:dyn-step-config}
Let $w$ be a digit string and $b\in\Nat$.  If
$\config w b\to_{\dynRules}v$, then $v=\config{w'}{b'}$ for some digit
string $w'$ and counter $b'$ with
\[
  \val(\config{w'}{b'})=\Bigl\lfloor\frac{3\,\val(\config w b)}2\Bigr\rfloor
\]
and
\[
  \begin{cases}
    b'=b+2, & \val(\config w b)\text{ even},\\
    b=b'+1, & \val(\config w b)\text{ odd}.
  \end{cases}
\]
\end{lemma}
\begin{proof}
By Lemma~\ref{lem:dyn-align} the rewrite is either the even rule
$w=w'\conc f\mapsto w'\conc0$ with $b'=b+2$, or the odd rule
$w=w'\conc t\mapsto w'\conc1$ with $b=b'+1$.  The value equations are
exactly Lemmas \ref{lem:dyn-even} and~\ref{lem:dyn-odd}.
\end{proof}

Auxiliary steps are even better behaved: they never change the shape of
a configuration.

\begin{lemma}[\texttt{auxStep\_config\_dest}]
\label{lem:aux-step-config}
Let $w$ be a digit string and $b\in\Nat$.  If
$\config w b\to_{\auxRules}v$, then $v=\config{w'}{b}$ for some digit
string $w'$.
\end{lemma}
\begin{proof}[Proof sketch]
Transport rules have left-hand sides consisting of two digit symbols
with no marker, so a redex inside $\config w b$ must lie entirely
inside the digit block $w$; boundary rules have left-hand side
$\lt d$ with $d$ a digit, so they rewrite the first digit of $w$.
Neither kind mentions $c$, so the counter $b$ is unchanged.
\end{proof}

Putting the two kinds of step together gives the main simulation theorem.
We write $\valOrbit{a_0}{n}$ for the $n$-th iterate of
$a\mapsto\lfloor3a/2\rfloor$ starting from $a_0$, and
$\counterZ{b_0}{a_0}{n}$ for the counter after $n$ macro steps when the
initial counter is $b_0$.

\begin{theorem}[\texttt{antihydraSRS\_simulates}]
\label{thm:faithful-simulation}
Let $w_0$ be a digit string and $b_0\in\Nat$, and let $a_0=\val(\config{w_0}{b_0})$.
For every string $C$ reachable from $\config{w_0}{b_0}$ by the full SRS
$\antihydraSRS$ there exist a digit string $w$, a counter $b$, and a
number $k\in\Nat$ of dynamic steps such that
\[
  C=\config w b,\qquad
  \val(\config w b)=\valOrbit{a_0}{k},\qquad
  b=\counterZ{b_0}{a_0}{k}.
\]
In particular, the value and counter after $k$ dynamic steps are
independent of the auxiliary rewriting strategy.
\end{theorem}
\begin{proof}
By induction on the reflexive-transitive closure of the rewrite
relation.  The base case $k=0$ is trivial.  For the inductive step,
use Lemma~\ref{lem:aux-step-config} for auxiliary steps (they keep $k$
and the counter unchanged and preserve the value by
Theorem~\ref{thm:aux-step}) and Lemma~\ref{lem:dyn-step-config} for a
dynamic step (it increments $k$ by one and updates value and counter by
exactly one macro step).
\end{proof}

Since the value orbit is deterministic, the number of even and odd
steps among the first $k$ dynamic steps is fixed.  Theorem~\ref{thm:faithful-simulation}
therefore implies the strategy-independence of the counter in a form
that refers directly to SRS derivations.

\begin{corollary}[\texttt{counter\_strategy\_independent}]
\label{cor:counter-strategy}
Let $w_0$ be a digit string and $b_0\in\Nat$, and let $a_0=\val(\config{w_0}{b_0})$.
For every configuration $C$ reachable from $\config{w_0}{b_0}$ there
exist $k$ and $j$ such that
\[
  C=\config w b\quad\text{and}\quad
  b=b_0+2k-3j,
\]
where $j$ is the number of odd steps among the first $k$ macro steps
from $a_0$.
\end{corollary}
\begin{proof}
Theorem~\ref{thm:faithful-simulation} gives $k$ with
$b=\counterZ{b_0}{a_0}{k}$.  Writing $e$ for the number of even steps
and $j$ for the number of odd steps among the first $k$ steps, we have
$e+j=k$ and $b=b_0+2e-j=b_0+2k-3j$.
\end{proof}

\subsubsection*{Comparison with the soundness layer.}
The file \texttt{AntihydraSRS.lean} proves local soundness: each rule
preserves or correctly updates value and counter.  The file
\texttt{AntihydraSRSSimulation.lean} lifts those local facts to a global
statement about arbitrary derivations.  This is exactly the step from
YAH's Claim~1 (per-rule correctness) to the general simulation lemma of
Theorem~3.17.

We can also exhibit the macro step as an explicit one-step rewrite on
configurations.

\begin{lemma}[\texttt{config\_even\_step}]
  \label{lem:config-even}
  For every value word $w$ and every counter $b$,
  \[
    \config{(w\conc f)}{b}
    \to_{\antihydraSRS}
    \config{(w\conc 0)}{b+2}.
  \]
\end{lemma}
\begin{proof}
  Apply the rule $f\gt\to0\gt cc$ in the context
  $\lt w$ on the left and $c^{b}$ on the right.  Formally,
  \[
    \config{(w\conc f)}{b}
    =
    (\lt w)\,(f\gt)\,c^{b},
    \qquad
    \config{(w\conc 0)}{b+2}
    =
    (\lt w)\,(0\gt cc)\,c^{b}.
  \]
\end{proof}

\begin{lemma}[\texttt{config\_odd\_step}]
  \label{lem:config-odd}
  For every value word $w$ and every counter $b$,
  \[
    \config{(w\conc t)}{b+1}
    \to_{\antihydraSRS}
    \config{(w\conc 1)}{b}.
  \]
\end{lemma}
\begin{proof}
  Apply the rule $t\gt c\to1\gt$ in the context
  $\lt w$ on the left and $c^{b}$ on the right.  The trailing $c$ of the
  rule consumes one counter symbol from the $c^{b+1}$ block, leaving
  $c^{b}$.
\end{proof}

\begin{lemma}[\texttt{dynRules\_odd\_needs\_counter}]
  \label{lem:odd-halt-redex}
  The string $t\gt$ is not the left-hand side of any dynamic rule.
  Consequently, a configuration of the form $u\,t\gt$ with no counter
  symbol is irreducible.
\end{lemma}
\begin{proof}
  The two dynamic left-hand sides are $f\gt$ and $t\gt c$; neither
  equals $t\gt$.
\end{proof}

The preceding lemma matches the macro halting condition.

\begin{theorem}[\texttt{nextMathState\_halt\_condition}]
  \label{thm:halt-bridge}
  For $a\ge2$, the BusyLean-free macro model halts with empty counter,
  $\nextMathState(a,0)=\none$, if and only if $a$ is odd.
\end{theorem}
\begin{proof}
  By definition $\nextMathState(a,0)$ returns $\none$ exactly on the
  odd branch of the conditional, which is taken precisely when $a$ is
  odd.  (The hypothesis $a\ge2$ guarantees that the outer guard
  $a/2\ge1$ holds.)
\end{proof}

Thus an odd value with an empty counter is a halting configuration both
for the SRS (Lemma~\ref{lem:odd-halt-redex}) and for the macro model
(Theorem~\ref{thm:halt-bridge}).

\subsection{Forward (constructive) simulation}
\label{subsec:lean-forward}

Theorem~\ref{thm:faithful-simulation} is the \emph{backward} half of a
simulation theorem: every SRS derivation stays inside the deterministic
orbit.  It is compatible, in principle, with the SRS making no progress
at all.  The file \texttt{AntihydraSRSForward.lean} supplies the missing
\emph{forward} half: an explicit, essentially deterministic rewriting
strategy that walks the orbit step by step.  Together the two halves give
an exact characterisation of the reachable configurations.

\subsubsection{Binary digit strings}

For the forward construction it is convenient to restrict attention to
value words that contain only binary digits.

\begin{definition}[\texttt{IsBinary}]
A list $w\in\List{\Sigma_A}$ is \emph{binary} if every symbol of $w$ is
$f$ or $t$.
\end{definition}

\begin{lemma}[\texttt{IsBinary.isDigits}]
Every binary string is a digit string.
\end{lemma}
\begin{proof}
A binary symbol is either $f$ or $t$, both of which are digits.
\end{proof}

\begin{lemma}[\texttt{IsBinary.snoc}]
If $w$ is binary and $s$ is $f$ or $t$, then $w\conc[s]$ is binary.
\end{lemma}
\begin{proof}
Immediate from the definition.
\end{proof}

\begin{lemma}[\texttt{IsBinary.exists\_snoc}]
\label{lem:binary-snoc}
Every nonempty binary string $w$ splits uniquely as $w=p\conc[s]$ where
$s\in\{f,t\}$ and $p$ is binary.
\end{lemma}
\begin{proof}
Take $s$ to be the last symbol and $p$ the prefix; the conditions are
satisfied by construction, and uniqueness follows from the fact that the
last symbol of a list is unique.
\end{proof}

On a binary string $w$ the value function has the simple form
\[
  \val(\lt w\gt)=\compFun_w(1),
\]
which is exactly the natural number whose binary expansion is
$1\conc w$ (with $f=0$ and $t=1$).  Hence the parity of the value is read
off the last symbol: $f$ gives an even value and $t$ gives an odd value.

\subsubsection{Lifting auxiliary derivations}

The carry sweep below is built purely from the auxiliary rules
$\auxRules=\auxA\cup\auxB$.  To make it a derivation of the full SRS we
need a trivial monotonicity lemma.

\begin{lemma}[\texttt{auxStep\_to\_antihydra}]
\label{lem:aux-to-full}
Every single $\auxRules$-rewrite step is an $\antihydraSRS$-rewrite step.
\end{lemma}
\begin{proof}
The full rule set is $\antihydraSRS=\dynRules\cup\auxRules$, so every
rule of $\auxRules$ is also a rule of $\antihydraSRS$.
\end{proof}

\begin{theorem}[\texttt{auxDeriv\_to\_antihydra}]
\label{thm:aux-deriv-to-full}
If $u\to^{*}_{\auxRules}v$, then $u\to^{*}_{\antihydraSRS}v$.
\end{theorem}
\begin{proof}
Reflexive-transitive closure is monotone with respect to inclusion of the
underlying relation.  Lemma~\ref{lem:aux-to-full} gives the inclusion of
single-step relations.
\end{proof}

\subsubsection{The carry sweep}

A dynamic firing replaces the rightmost binary digit of $w$ by a single
ternary digit ($f\mapsto0$, $t\mapsto1$).  Before the next dynamic step
can fire, that ternary digit must be carried leftward through the binary
block until it reaches the left boundary $\lt$, where it is resolved back
into binary digits.  The transport rules $\auxA$ perform one leftward
swap at a time.

\begin{lemma}[\texttt{auxA\_swap}]
\label{lem:auxa-swap}
Let $s\in\{f,t\}$ be a binary symbol and $d_k\in\{0,1,2\}$ a ternary
symbol.  Then there exist a ternary symbol $d_j$ and a binary symbol
$s_j$ such that $s\,d_k\to d_j\,s_j$ is a rule of $\auxA$ and, for all
$x\in\Nat$,
\[
  \symfun_{d_k}\bigl(\symfun_s(x)\bigr)
  =
  \symfun_{s_j}\bigl(\symfun_{d_j}(x)\bigr).
\]
\end{lemma}
\begin{proof}
There are six cases, corresponding to the six transport rules in
\eqref{eq:lean-transport}.  Each is a direct computation.  For instance,
for $t1\to2f$ we have
$\symfun_1(\symfun_t(x))=3(2x+1)+1=6x+4=2(3x+2)=\symfun_f(\symfun_2(x))$.
The other five cases are analogous.
\end{proof}

\begin{theorem}[\texttt{carry\_sweep}]
\label{thm:carry-sweep}
Let $p$ be a binary string and $d_k\in\{0,1,2\}$.  For every suffix
$\textit{rest}$ there exists a binary string $q$ such that
\[
  \val(\lt q\gt)=\symfun_{d_k}\bigl(\val(\lt p\gt)\bigr)
\]
and
\[
  \lt\,p\,d_k\,\textit{rest}
  \to^{*}_{\auxRules}
  \lt\,q\,\textit{rest}.
\]
\end{theorem}
\begin{proof}
By reverse induction on $p$.  If $p=\varepsilon$, the string is
$\lt\,d_k\,\textit{rest}$ and a single boundary rule resolves it:
$\lt 0\to\lt t$, $\lt 1\to\lt ff$, or $\lt 2\to\lt ft$.  The resulting $q$
is $[t]$, $[f,f]$, or $[f,t]$, and the value equation is immediate.

If $p=p'\conc[s]$ with $s\in\{f,t\}$, Lemma~\ref{lem:auxa-swap} gives
$d_j$ and $s_j$ with $s\,d_k\to d_j\,s_j\in\auxA$ and equal composite
functions.  One transport step yields
\[
  \lt\,p'\,s\,d_k\,\textit{rest}
  \to_{\auxA}
  \lt\,p'\,d_j\,s_j\,\textit{rest}.
\]
The induction hypothesis applied to $p'$, $d_j$, and
$\textit{rest}'=s_j\,\textit{rest}$ gives a binary $q'$ with
$\val(\lt q'\gt)=\symfun_{d_j}(\val(\lt p'\gt))$ and
$\lt\,p'\,d_j\,s_j\,\textit{rest}\to^{*}_{\auxRules}\lt\,q'\,s_j\,\textit{rest}$.
Set $q=q'\conc[s_j]$.  Then $q$ is binary by
Lemma~\ref{lem:binary-snoc}, and
\[
  \begin{aligned}
  \val(\lt q\gt)
  &=2\,\val(\lt q'\gt)+\delta\\
  &=\symfun_{s_j}\bigl(\symfun_{d_j}(\val(\lt p'\gt))\bigr)\\
  &=\symfun_{d_k}\bigl(\symfun_s(\val(\lt p'\gt))\bigr)\\
  &=\symfun_{d_k}\bigl(\val(\lt p\gt)\bigr).
  \end{aligned}
\]
where $\delta=0$ if $s_j=f$ and $\delta=1$ if $s_j=t$.  Concatenating the
single transport step with the induction hypothesis gives the required
derivation.
\end{proof}

\subsubsection{The macro orbit as an option iteration}

To state the forward theorem compactly, the Lean file defines an
iteration of the partial macro step.

\begin{definition}[\texttt{macroIter}]
Define $\macroIter:\Nat\to(\Nat\times\Nat)\to\Option(\Nat\times\Nat)$
by
\[
  \macroIter(0,s)=\some(s),
  \qquad
  \macroIter(n+1,s)=\macroIter(n,s)\bind\macroStep,
\]
where $u\bind f$ is $\none$ if $u=\none$ and $f(a,b)$ if $u=\some(a,b)$.
\end{definition}

Thus $\macroIter(n,(a_0,b_0))=\some(a_n,b_n)$ means precisely that the
macro model has not halted within the first $n$ steps and that
$(a_n,b_n)$ is the state after $n$ steps.

\subsubsection{The forward simulation theorem}

We can now prove that the SRS does not merely stay inside the orbit, but
actually walks it.

\begin{theorem}[\texttt{antihydraSRS\_realizes\_orbit}]
\label{thm:forward-simulation}
Let $w_0$ be a binary string and $b_0\in\Nat$, and assume
$a_0=\val(\lt w_0\gt)\ge2$.  For every $n,a_n,b_n\in\Nat$ with
\[
  \macroIter(n,(a_0,b_0))=\some(a_n,b_n),
\]
there exists a binary string $w_n$ such that
\[
  \val(\lt w_n\gt)=a_n,\qquad \val(\lt w_n\gt)\ge2,
\]
and
\[
  \config{w_0}{b_0}
  \to^{*}_{\antihydraSRS}
  \config{w_n}{b_n}.
\]
\end{theorem}
\begin{proof}
Induction on $n$.  The base case $n=0$ is trivial: take $w_0=w_n$.

For the inductive step, suppose
$\macroIter(n+1,(a_0,b_0))=\some(a_{n+1},b_{n+1})$.  Unfolding the
iteration gives $(a_n,b_n)$ with
$\macroIter(n,(a_0,b_0))=\some(a_n,b_n)$ and
$\macroStep(a_n,b_n)=\some(a_{n+1},b_{n+1})$.  By the induction
hypothesis there is a binary $w_n$ with $\val(\lt w_n\gt)=a_n\ge2$ and
$\config{w_0}{b_0}\to^{*}_{\antihydraSRS}\config{w_n}{b_n}$.

Because $a_n\ge2$, the word $w_n$ is nonempty, so by
Lemma~\ref{lem:binary-snoc} it splits as $w_n=p\conc[s]$ with
$s\in\{f,t\}$.  Consider the two parity cases.

\paragraph{Even case ($s=f$).}  Here
$\val(\lt w_n\gt)=2\,\val(\lt p\gt)$, so $a_n=2\,\val(\lt p\gt)$ and
\[
  \macroStep(a_n,b_n)=(3\,\val(\lt p\gt),\,b_n+2).
\]
Lemma~\ref{lem:config-even} gives the single dynamic step
$\config{(p\conc[f])}{b_n}\to_{\antihydraSRS}\config{(p\conc[0])}{b_n+2}$.
Theorem~\ref{thm:carry-sweep}, with $d_k=0$ and
$\textit{rest}=\gt\,c^{b_n+2}$, yields a binary $q$ with
$\val(\lt q\gt)=3\,\val(\lt p\gt)$ and
$\lt\,p\,0\,\gt\,c^{b_n+2}\to^{*}_{\auxRules}\lt\,q\,\gt\,c^{b_n+2}$.
By Theorem~\ref{thm:aux-deriv-to-full} this is also a derivation of
$\antihydraSRS$, so
$\config{(p\conc[0])}{b_n+2}\to^{*}_{\antihydraSRS}\config{q}{b_n+2}$.
Set $w_{n+1}=q$; then $(a_{n+1},b_{n+1})=(\val(\lt q\gt),b_n+2)$ as
required.

\paragraph{Odd case ($s=t$).}  Here
$\val(\lt w_n\gt)=2\,\val(\lt p\gt)+1$, so $a_n=2\,\val(\lt p\gt)+1$.
Since $\macroStep(a_n,b_n)$ is defined, the counter is positive; write
$b_n=b'+1$.  Then
\[
  \macroStep(a_n,b_n)=(3\,\val(\lt p\gt)+1,\,b').
\]
Lemma~\ref{lem:config-odd} gives
$\config{(p\conc[t])}{b'+1}\to_{\antihydraSRS}\config{(p\conc[1])}{b'}$,
and Theorem~\ref{thm:carry-sweep} with $d_k=1$ carries the digit $1$ to a
binary $q$ with $\val(\lt q\gt)=3\,\val(\lt p\gt)+1$.  Again set
$w_{n+1}=q$.

In both cases we have extended the derivation from $\config{w_0}{b_0}$ to
$\config{w_{n+1}}{b_{n+1}}$, completing the induction.
\end{proof}

The value produced by $\macroIter$ is exactly the deterministic value
orbit.

\begin{lemma}[\texttt{macroIter\_fst\_eq\_valOrbit}]
\label{lem:macroiter-valorbit}
For all $a_0,b_0,n,a_n,b_n$, if
\[
  \macroIter(n,(a_0,b_0))=\some(a_n,b_n),
\]
then $a_n=\valOrbit{a_0}{n}$.
\end{lemma}
\begin{proof}
By induction on $n$, using the fact that $\macroStep$ applies the map
$a\mapsto\lfloor3a/2\rfloor$ on both the even and the odd branch
(Lemmas \ref{lem:dyn-even} and~\ref{lem:dyn-odd}).
\end{proof}

Repackaging Theorem~\ref{thm:forward-simulation} against the value orbit
gives the following direct converse of
Theorem~\ref{thm:faithful-simulation}.

\begin{theorem}[\texttt{antihydraSRS\_realizes\_valOrbit}]
\label{thm:forward-valorbit}
Let $w_0$ be binary with $a_0=\val(\lt w_0\gt)\ge2$, and let
$\macroIter(n,(a_0,b_0))=\some(a_n,b_n)$.  Then there exists a binary
$w_n$ such that
\[
  \val(\lt w_n\gt)=\valOrbit{a_0}{n}
  \qquad\text{and}\qquad
  \config{w_0}{b_0}
  \to^{*}_{\antihydraSRS}
  \config{w_n}{b_n}.
\]
\end{theorem}
\begin{proof}
Combine Theorem~\ref{thm:forward-simulation} with
Lemma~\ref{lem:macroiter-valorbit}.
\end{proof}

Finally, the backward and forward theorems together give an exact
characterisation of reachability.

\begin{theorem}[\texttt{antihydraSRS\_orbit\_exact}]
\label{thm:orbit-exact}
Let $w_0$ be binary with $a_0=\val(\lt w_0\gt)\ge2$ and let
$\macroIter(n,(a_0,b_0))=\some(a_n,b_n)$.  Then:
\begin{enumerate}
\item There exists a binary $w_n$ with
  $\val(\lt w_n\gt)=\valOrbit{a_0}{n}$ and
  $\config{w_0}{b_0}\to^{*}_{\antihydraSRS}\config{w_n}{b_n}$.
\item Conversely, for every configuration $C$ reachable from
  $\config{w_0}{b_0}$ there exist $k$, $w$, and $b$ such that
  $C=\config{w}{b}$ and $\val(\lt w\gt)=\valOrbit{a_0}{k}$.
\end{enumerate}
In other words, the configurations reachable from
$\config{w_0}{b_0}$ are exactly the orbit configurations.
\end{theorem}
\begin{proof}
Part~(1) is Theorem~\ref{thm:forward-valorbit}; part~(2) is
Theorem~\ref{thm:faithful-simulation}.
\end{proof}

\begin{corollary}[\texttt{antihydraSRS\_realizes\_orbit\_from\_eight}]
\label{cor:from-eight}
For every $b_0,n,a_n,b_n$ with
$\macroIter(n,(8,b_0))=\some(a_n,b_n)$, there exists a binary $w_n$ such
that
$\val(\lt w_n\gt)=\valOrbit{8}{n}$ and
$\config{[f,f,f]}{b_0}\to^{*}_{\antihydraSRS}\config{w_n}{b_n}$.
\end{corollary}
\begin{proof}
The word $[f,f,f]$ is binary and $\val(\lt fff\gt)=8$, so
Theorem~\ref{thm:forward-valorbit} applies.
\end{proof}

\subsubsection*{Comparison with the soundness and backward layers}
The chain of files is now complete.  \texttt{AntihydraSRS.lean} proves
local soundness (each rule preserves or correctly updates value and
counter); \texttt{AntihydraSRSSimulation.lean} lifts this to the global
backward statement that arbitrary derivations are faithful to the orbit;
and \texttt{AntihydraSRSForward.lean} supplies the constructive forward
strategy that realises the orbit.  Theorem~\ref{thm:orbit-exact} is the
exact analogue, for Antihydra, of YAH's Theorem~3.17 in full: Claims~1
and~2 give the backward simulation, and Claim~3 gives the forward
realisation.

\subsection{Bridge to the Turing machine}
\label{subsec:tm-bridge}

The file \texttt{AntihydraMachine.lean} contains the BusyLean proof
that the macro model simulates the actual Antihydra Turing machine.  The
main ingredients are the six macro theorems of
Section~\ref{subsec:macro}, which are lifted from finite padded
configurations $P(b,m,n,p)$ to stream configurations $SP(b,m,n)$ and
assembled into the simulation theorem
\texttt{stm\_simulates\_math}.  This theorem says that for every
$a\ge2$ and every $b$, there is a positive number of TM steps $k$ such
that the machine from $SP(b,a,0)$ reaches $SP(b',a',0)$ if
$\nextMathState(a,b)=\some(a',b')$, and halts if
$\nextMathState(a,b)=\none$.

The halting predicate $\mathHalts$ is defined inductively from
$\nextMathState$ in the usual way, and the crucial bridge
(\texttt{mathHalts\_iff\_Aiter\_8\_2} in
\texttt{AntihydraSRSaxiom.lean}) states that
\begin{equation}
  \label{eq:halts-iter}
  \mathHalts(8,2)
  \;\iff\;
  \exists i\;\bigl(A^{i}(8,2).1\text{ is odd}\;\wedge\;A^{i}(8,2).1/2\ge1\;\wedge\;A^{i}(8,2).2=0\bigr),
\end{equation}
where $A^{i}$ is the $i$-th iterate of the total macro map $A$.  In
\texttt{AntihydraMachine.lean} this equivalence is proved directly from
the definitions of $\nextMathState$ and $\mathHalts$, and then combined
with \texttt{stm\_simulates\_math} to give the main TM-level result
\texttt{antihydra\_halts\_iff}: the Antihydra TM halts if and only if
the right-hand side of \eqref{eq:halts-iter} holds.

Composing this with Theorem~\ref{thm:srs-represents-macro} and
Theorem~\ref{thm:halt-bridge}, we obtain the intended interpretation of
the SRS:

\begin{corollary}
  \label{cor:srs-tm}
  The Antihydra SRS is non-terminating from the initial configuration
  $\lt fff\gt$ if and only if the Antihydra Turing machine does not
  halt.
\end{corollary}

\begin{proof}[Proof sketch]
  By Lemma~\ref{lem:initial-config} below, $\lt fff\gt$ represents
  $(a,b)=(8,0)$.  The mixed-base macro step $\macroStep$ and the
  $m$-coordinate partial step $\nextMathState$ have the same halting
  condition (Theorem~\ref{thm:halt-bridge}), and the total map $A$ is
  the totalisation of $\nextMathState$.  The BusyLean proof identifies
  non-halting of the TM with the statement that no iterate of $A$ from
  the appropriate starting point reaches an odd value with empty counter.
  Hence a non-terminating SRS derivation corresponds exactly to an
  infinite Antihydra orbit, and a halting SRS configuration corresponds
  exactly to a halting TM configuration.
\end{proof}

\paragraph{Relation to the general YAH simulation lemma.}
The proof template used above is the Antihydra analogue of
Theorem~3.17 of Yolcu--Aaronson--Heule~\cite[pp.~20--21]{YAH21}: one
introduces a mixed-base SRS whose auxiliary rules preserve the
represented integer value while the dynamic rules implement the iterated
map, then proves that the SRS is non-terminating exactly when the
original system does not halt.  In the Antihydra case the dynamic rules
realise the lower $3/2$-map with a unary counter, whereas YAH's original
system realises the Collatz map~$T$.  The same strategy-independence,
value-preservation, and halting-characterization lemmas apply whenever
the macro dynamics can be written as a piecewise-affine digit update on a
mixed-base representation.  The broader research program, recorded in
Report~B3 (\S8) and in \texttt{3-forward-strategy.html}
(\S5)~\cite{ReportB3,ForwardStrategy}, is to extract this pattern into a
single general simulation lemma for the full YAH rewriting framework,
prove it once in Lean, and use it as a verified front-end for all
certificate-producing termination proofs of Collatz-like systems.  The
formalization in this section is a worked instance of that program.

\subsection{Initial configuration}

\begin{lemma}[\texttt{initial\_config}]
  \label{lem:initial-config}
  The configuration $\config{[f,f,f]}{0}=\lt fff\gt$ has value $8$ and
  counter $0$.
\end{lemma}
\begin{proof}
  Compute the left fold:
  $0\xrightarrow{\lt}1\xrightarrow{f}2\xrightarrow{f}4\xrightarrow{f}8\xrightarrow{\gt}8$.
  There are no counter symbols.
\end{proof}

\section{FAR/CPS window saturation}
\label{sec:far}

A \emph{window} of width $w$ over a string $s$ is a length-$w$ factor of
the $(w-1)$-padded string $\lpad^{\,w-1}s\,\rpad^{\,w-1}$, where \lpad
and \rpad are fresh left/right padding symbols.  The FAR/CPS method
searches for a regular (finite-window) superset of all reachable
configurations: a set $W$ of width-$w$ windows that contains every window
of the initial configuration and is closed under the rewriting rules.

\paragraph{Saturation procedure.}
Starting from $W_0=\grams{\lt fff\gt}{w}$, repeatedly apply every rule
$\ell\to r$ in every left context of length $w-1$ and every right
context of length $w-1$ whose windows are already known to be reachable,
and add all windows of the resulting string to the set.  More
formally, if $u\ell v$ is such that every length-$w$ window of
$u\ell v$ lies in the current $W$, then every length-$w$ window of
$urv$ is added to $W$.  The process is repeated until a fixpoint is
reached.  The implementation is in \texttt{antihydra\_far.py}, functions
\texttt{grams}, \texttt{left\_contexts}, \texttt{right\_contexts}, and
\texttt{saturate}.

When the fixpoint is computed, two properties are checked:
\begin{itemize}
\item \textbf{Control:} the window $2\gt$ (ternary digit $2$ adjacent
to the boundary) is absent from $W$.  This is a genuine invariant:
Antihydra's even dynamic rule appends ternary digit $0$ and the odd
rule appends ternary digit $1$, so a ternary ending $2$ can never
appear.
\item \textbf{Target:} no string ending in $t\gt$ (an odd value with
empty counter, i.e.\ a halting configuration) has all its windows in
$W$.  If such a string exists, it is a \emph{phantom}: the abstraction
admits a halting configuration that is not actually reachable.
\end{itemize}

Table~\ref{tab:far} shows the results for widths $2$--$6$.

\begin{table}[ht]
\centering
\caption{FAR/CPS window saturation for Antihydra.}
\label{tab:far}
\begin{tabular}{rrrrlll}
\toprule
$w$ & $|W|$ & passes & time & control & target & phantom \\
\midrule
2 & 39 & 7 & $<0.01$s & yes & no & $\lt t\gt$ \\
3 & 185 & 11 & $0.01$s & yes & no & $\lt ft\gt$ \\
4 & 889 & 18 & $0.35$s & yes & no & $\lt ftt\gt$ \\
5 & 4333 & 29 & $16$s & yes & no & $\lt fttt\gt$ \\
6 & 21187 & 37 & $542$s ($\approx$9\,min) & yes & no & $\lt tftt\gt$ \\
\bottomrule
\end{tabular}
\end{table}

\paragraph{The phantom family.}
At every width the target fails: the abstraction admits a halting
configuration that is not actually reachable.  For widths $2$--$5$ the
minimal such phantom is
\[
  \lt f\,t^{\,w-2}\gt
\]
(with the convention that for $w=2$ this is $\lt t\gt$).  Its value is
\[
  \val(\lt f\,t^{\,w-2}\gt)=3\cdot2^{\,w-2}-1,
\]
an odd number whose binary expansion ends in a run of $w-2$ ones.  Such
a value commits to $w-2$ consecutive odd macro steps, each decreasing
the counter by $1$, but the abstraction sees an empty counter at the
boundary and cannot compare the length of the odd run with the amount
of counter available.  These are the \emph{$-1$-shadow} configurations
of Deliverable~4: integers locally indistinguishable from the
(negative) fixed point $-1$, whose odd-run length is the $2$-adic budget
$\nu_2(a+1)$.  At width $w=6$ a shorter phantom, $\lt tftt\gt$,
appears; it is admitted by the width-$6$ window language while the
previous phantom $\lt ftttt\gt$ is excluded.  Thus each increment of
$w$ kills some phantoms and admits new ones, and no finite window width
can exclude all halting configurations.

\section{Pumping obstruction to regular invariants}
\label{sec:obstruction}

A FAR-style non-halting proof exhibits a regular set $S$ of configurations
that (i) contains every reachable configuration, (ii) is closed under one
rewriting step, and (iii) contains no halting configuration.  The file
\texttt{AntihydraSRSObstruction.lean} proves a hard limitation of such
proofs for Antihydra when the counter is stored as a contiguous unary
block: if $S$ is recognized with $p$ automaton states, then the future
dips of the counter are bounded by $p$.

Let $(a_n,b_n)$ be the state of the deterministic macro orbit after $n$
steps, starting from a binary configuration $\config{w_0}{b_0}$ with
$a_0=\val(\lt w_0\gt)\ge2$.  Define the \emph{future-dip sequence}
\begin{equation}
  \label{eq:future-dip}
  D(n)=b_n-\min_{m\ge n}b_m.
\end{equation}
Thus $D(n)$ measures how far the counter falls from its current level
before recovering.  The main result is:

\begin{proposition}[Pumping obstruction]
\label{prop:pumping-obstruction}
Let $S$ be a set of configurations satisfying \textup{(i)--(iii)} above
and suppose $S$ admits unary-counter pumping below $p$: whenever
$\config{w}{b}\in S$ with $b\ge p$, there exists $b'<p$ such that
$\config{w}{b'}\in S$.  Then for every $n$ with $b_n\ge p$,
\[
  D(n)<p.
\]
Equivalently, $b_n<b_m+p$ for all $m\ge n$ whenever $b_n\ge p$.
\end{proposition}

Before giving the proof, we isolate the necessary orbit arithmetic and the
notion of a halting configuration.

\subsection{Orbit arithmetic}

The following lemmas are proved in
\texttt{AntihydraSRSObstruction.lean}.  They say that the value orbit is
additive in the step index, that shifting the initial counter shifts the
counter orbit by the same amount, and that the partial iterate
$\macroIter$ agrees with the deterministic orbits while it has not halted.

\begin{lemma}[\texttt{valOrbit\_add}]
\label{lem:valOrbit-add}
For all $a_0,n,j\in\Nat$,
\[
  \valOrbit{a_0}{n+j}=\valOrbit{\bigl(\valOrbit{a_0}{n}\bigr)}{j}.
\]
\end{lemma}
\begin{proof}
By induction on $j$, using the definition
$\valOrbit{a}{j+1}=\bigl\lfloor\frac32\valOrbit{a}{j}\bigr\rfloor$.
\end{proof}

\begin{lemma}[\texttt{counterZ\_shift}]
\label{lem:counterZ-shift}
For all integers $b,b'$ and all $a_0,j\in\Nat$,
\[
  \counterZ{b'}{a_0}{j}=\counterZ{b}{a_0}{j}+(b'-b).
\]
\end{lemma}
\begin{proof}
Both sides satisfy the same recurrence in $j$ with the same initial value,
because counter increments depend only on the value orbit, not on the
absolute counter level.
\end{proof}

\begin{lemma}[\texttt{counterZ\_add}]
\label{lem:counterZ-add}
For all $b_0\in\ZZ$, $a_0,n,j\in\Nat$,
\[
  \counterZ{b_0}{a_0}{n+j}
  =
  \counterZ{\bigl(\counterZ{b_0}{a_0}{n}\bigr)}{\bigl(\valOrbit{a_0}{n}\bigr)}{j}.
\]
\end{lemma}
\begin{proof}
By induction on $j$, using Lemma~\ref{lem:valOrbit-add} and the recurrence
for $\counterZ{}{}{}$.
\end{proof}

\begin{lemma}[\texttt{macroIter\_eq\_orbit}]
\label{lem:macroiter-eq-orbit}
If $\macroIter(n,(a_0,b_0))=\some(a_n,b_n)$, then
$a_n=\valOrbit{a_0}{n}$ and $b_n=\counterZ{b_0}{a_0}{n}$.
\end{lemma}
\begin{proof}
By induction on $n$, using the fact that $\macroStep$ applies
$a\mapsto\lfloor3a/2\rfloor$ on both parity branches and updates the
counter by $+2$ or $-1$ accordingly.
\end{proof}

\begin{lemma}[\texttt{valOrbit\_ge\_two}]
\label{lem:valOrbit-ge-two}
If $a_0\ge2$ then $\valOrbit{a_0}{n}\ge2$ for every $n$.
\end{lemma}
\begin{proof}
The map $a\mapsto\lfloor3a/2\rfloor$ preserves the property $a\ge2$.
\end{proof}

\begin{lemma}[\texttt{macroStep\_eq\_none\_iff}]
\label{lem:macroStep-none}
$\macroStep(a,b)=\none$ if and only if $a$ is odd and $b=0$.
\end{lemma}
\begin{proof}
Immediate from the definition of $\macroStep$.
\end{proof}

\begin{lemma}[\texttt{exists\_first\_none}]
\label{lem:first-none}
If $\macroIter(N,(a_0,b_0))=\none$, then there exist $j,a_j,b_j$ such
that $\macroIter(j,(a_0,b_0))=\some(a_j,b_j)$ and
$\macroIter(j+1,(a_0,b_0))=\none$.
\end{lemma}
\begin{proof}
By induction on $N$.
\end{proof}

\begin{lemma}[\texttt{mem\_of\_reflTransGen}]
\label{lem:mem-reflTransGen}
If $S$ is closed under one rewriting step and $u\in S$, then every
configuration $v$ with $u\to^{*}_{\antihydraSRS}v$ lies in $S$.
\end{lemma}
\begin{proof}
By induction on the reflexive-transitive closure.
\end{proof}

\subsection{Halting configurations}

\begin{definition}[\texttt{IsHalting}]
\label{def:halting}
A configuration is \emph{halting} if it has the form $\config{w}{0}$
where $w$ is binary and $\val(\lt w\gt)$ is odd.
\end{definition}

Such a configuration is irreducible: the only applicable dynamic rule for
an odd value is $t\gt c\to1\gt$, which requires a counter symbol $c$; the
counter is empty, so no rule fires.  By Lemma~\ref{lem:macroStep-none} this
is exactly the macro halting condition.

\subsection{Proof of the pumping obstruction}

\begin{proof}[Proof of Proposition~\ref{prop:pumping-obstruction}]
Fix $n$ with $b_n\ge p$ and any $m\ge n$.  By
Theorem~\ref{thm:forward-simulation} there is a binary word $W_n$ with
$\val(\lt W_n\gt)=a_n$ and a derivation
\[
  \config{w_0}{b_0}\to^{*}_{\antihydraSRS}\config{W_n}{b_n}.
\]
Since $S$ contains all reachable configurations,
$\config{W_n}{b_n}\in S$.  The pumping hypothesis gives $b'<p$ such that
$\config{W_n}{b'}\in S$.

Suppose, for a contradiction, that $b_n\ge b_m+p$.  Consider the partial
iteration from $(a_n,b')$ for $m-n$ steps.  We claim it halts, i.e.
\[
  \macroIter(m-n,(a_n,b'))=\none.
\]
If not, suppose it returns $\some(a'',b'')$.  By
Lemma~\ref{lem:macroiter-eq-orbit},
\[
  b''=\counterZ{b'}{a_n}{m-n}.
\]
Using Lemma~\ref{lem:counterZ-shift} and Lemma~\ref{lem:counterZ-add},
\[
  \counterZ{b'}{a_n}{m-n}
  =\counterZ{b_n}{a_n}{m-n}+(b'-b_n)
  =b_m+b'-b_n,
\]
which is negative because $b_n\ge b_m+p$ and $b'<p$.  This is impossible
for a natural-number counter, so the iteration must indeed halt.

By Lemma~\ref{lem:first-none} there is a first halting transition from
$(a_n,b')$: some $j$ with
\[
  \macroIter(j,(a_n,b'))=\some(a_j,b_j)
\]
and
\[
  \macroIter(j+1,(a_n,b'))=\none.
\]
Lemma~\ref{lem:macroStep-none} gives that $a_j$ is odd and $b_j=0$.
Apply
Theorem~\ref{thm:forward-simulation} again, this time starting from
$\config{W_n}{b'}$: there is a binary word $W_j$ with
$\val(\lt W_j\gt)=a_j$ (hence odd) and
\[
  \config{W_n}{b'}\to^{*}_{\antihydraSRS}\config{W_j}{0}.
\]
By Lemma~\ref{lem:mem-reflTransGen} and closure of $S$, the halting
configuration $\config{W_j}{0}$ lies in $S$, contradicting assumption~(iii).

Therefore $b_n<b_m+p$ for every $m\ge n$, i.e. $D(n)<p$.
\end{proof}

\subsection{Regular invariants and the corollary}

The hypotheses of Proposition~\ref{prop:pumping-obstruction} are bundled
in \texttt{AntihydraSRSObstruction.lean} into a structure.

\begin{definition}[\texttt{RegularInvariant}]
\label{def:regular-invariant}
A \emph{regular invariant} with $p$ states for the orbit from
$\config{w_0}{b_0}$ is a set $S$ of configurations such that:
\begin{enumerate}
\item[\textup{(reach)}] every configuration reachable from
  $\config{w_0}{b_0}$ lies in $S$;
\item[\textup{(closed)}] if $u\in S$ and $u\to_{\antihydraSRS}v$, then
  $v\in S$;
\item[\textup{(noHalt)}] $S$ contains no halting configuration;
\item[\textup{(pump)}] if $\config{w}{b}\in S$ with $b\ge p$, then there
  exists $b'<p$ with $\config{w}{b'}\in S$.
\end{enumerate}
\end{definition}

The \emph{orbit counter} is defined by
\[
  \orbitCounter{w_0}{b_0}{n}=(\macroIter(n,(\val(\lt w_0\gt),b_0))).\getD((0,0)).2,
\]
which reads off $b_n$ while the orbit has not halted.  With this notation,
Proposition~\ref{prop:pumping-obstruction} becomes:

\begin{theorem}[\texttt{orbitDip\_lt}]
\label{thm:orbit-dip}
Let $S$ be a regular invariant with $p$ states for the orbit from
$\config{w_0}{b_0}$.  Then for all $n,m$ with $n\le m$ and
$\orbitCounter{w_0}{b_0}{n}\ge p$,
\[
  \orbitCounter{w_0}{b_0}{n}
  <
  \orbitCounter{w_0}{b_0}{m}+p.
\]
\end{theorem}
\begin{proof}
Use Lemma~\ref{lem:macroiter-eq-orbit} to replace
$\orbitCounter{w_0}{b_0}{n}$ by $b_n$ and
$\orbitCounter{w_0}{b_0}{m}$ by $b_m$, then apply
Proposition~\ref{prop:pumping-obstruction}.
\end{proof}

In the language of future dips:

\begin{theorem}[\texttt{futureDip\_lt\_card}]
\label{thm:future-dip-lt}
Let $S$ be a regular invariant with $p$ states.  Then for every $n$ with
$\orbitCounter{w_0}{b_0}{n}\ge p$,
\[
  \futureDip(\orbitCounter{w_0}{b_0}{\cdot},n)<p.
\]
\end{theorem}

\begin{corollary}[\texttt{no\_regular\_invariant\_of\_unbounded\_dips}]
\label{cor:no-regular-invariant}
If the future dips of the orbit counter are unbounded, i.e.
\[
  \forall q\;\exists n\quad
  \bigl(\orbitCounter{w_0}{b_0}{n}\ge q\;\wedge\;
  \futureDip(\orbitCounter{w_0}{b_0}{\cdot},n)\ge q\bigr),
\]
then no regular invariant of any size exists for the orbit from
$\config{w_0}{b_0}$.
\end{corollary}
\begin{proof}
A $p$-state invariant would force
\[
  \futureDip(\orbitCounter{w_0}{b_0}{\cdot},n)<p
\]
whenever $\orbitCounter{w_0}{b_0}{n}\ge p$
(Theorem~\ref{thm:future-dip-lt}),
contradicting unboundedness at $q=p$.
\end{proof}

\subsection{The pumping property from a finite automaton}

The abstract pumping property~(pump) in
Definition~\ref{def:regular-invariant} is exactly what a DFA recognizer
with a contiguous unary counter block provides.  We make this literal.

\begin{lemma}[\texttt{exists\_iterate\_lt\_card}]
\label{lem:iterate-pigeonhole}
Let $\sigma$ be a finite type, $g:\sigma\to\sigma$, $q\in\sigma$, and
$n\ge|\sigma|$.  Then there exists $k<|\sigma|$ such that
$g^{[k]}(q)=g^{[n]}(q)$.
\end{lemma}
\begin{proof}
The sequence $g^{[0]}(q),g^{[1]}(q),\dots,g^{[n]}(q)$ has $n+1>|
sigma|$ entries, so two of them coincide by the pigeonhole principle.  The
periodic behaviour of iterates then reduces the larger index to a smaller
one.
\end{proof}

\begin{lemma}[\texttt{dfa\_pump}]
\label{lem:dfa-pump}
Let $M$ be a DFA over $\Sigma_A$ with state set $\sigma$ and let
$p=|\sigma|$.  If $M$ accepts $\config{w}{b}$ with $b\ge p$, then $M$
also accepts $\config{w}{b'}$ for some $b'<p$.
\end{lemma}
\begin{proof}
After reading the value prefix $\lt w\gt$, the DFA enters some state
$q$.  Reading the counter block $c^{b}$ iterates the $c$-transition
$b$ times.  Since $b\ge p=|\sigma|$,
Lemma~\ref{lem:iterate-pigeonhole} gives $k<p$ with the same final state
as $b$, so $M$ accepts $\config{w}{k}$.
\end{proof}

\begin{theorem}[\texttt{regularInvariant\_of\_dfa}]
\label{thm:regularInvariant-of-dfa}
Let $M$ be a DFA with state set $\sigma$ whose accepted language $L(M)$
contains every configuration reachable from $\config{w_0}{b_0}$, is closed
under one $\antihydraSRS$-step, and contains no halting configuration.
Then $L(M)$ is a regular invariant with $p=|\sigma|$ states.
\end{theorem}
\begin{proof}
Hypotheses~(i)--(iii) are exactly reachability, closure, and no-halt.
Lemma~\ref{lem:dfa-pump} supplies property~(pump).
\end{proof}

\begin{corollary}[\texttt{no\_dfa\_invariant\_of\_unbounded\_dips}]
\label{cor:no-dfa-invariant}
If the future dips of the orbit counter from $\config{w_0}{b_0}$ are
unbounded, then no DFA recognizes a closed, reachable-containing,
halting-free set of configurations.  In particular, no FAR-style
non-halting proof of Antihydra over unary-counter encodings exists at any
number of states.
\end{corollary}
\begin{proof}
Combine Theorem~\ref{thm:regularInvariant-of-dfa} with
Corollary~\ref{cor:no-regular-invariant}.
\end{proof}

\subsection{Relation to the FAR/CPS search}

The window-saturation experiments of Section~\ref{sec:far} produce, for
each width $w$, a regular over-approximation of the reachable
configurations.  The obstruction of this section says that any genuine
regular invariant for Antihydra must bound future dips by the number of
automaton states.  Since the Antihydra counter performs a random walk with
positive drift, future dips are believed to be unbounded; if that is true,
then no finite automaton can witness non-halting, and the FAR/CPS method is
fundamentally incomplete for this machine.  The phantom configurations
observed in Table~\ref{tab:far} are the concrete symptom: the window
abstraction is too weak to remember the counter budget needed to survive
long odd runs.

It is worth noting the logical strength of the conclusion.  The statement
``no finite automaton can witness non-halting'' is weaker than the
statement ``Antihydra does not halt.''  The pumping obstruction gives
$\text{FAR witness}\Rightarrow\text{dips bounded}$, so unbounded dips
yields two consequences: (1)~the machine does not halt, because halting
would cap every future dip by the maximum counter value ever reached;
and (2)~no FAR-style certificate exists.  Thus unbounded dips is a
strictly stronger result than either conclusion alone; the FAR
incompleteness theorem is the additional meta-mathematical information
that one natural class of certificates is too weak to establish the
true non-halting behaviour.

\begin{thebibliography}{9}
\bibitem{YAH21}
Emre Yolcu, Scott Aaronson, and Marijn J.~H. Heule,
\textit{An Automated Approach to the Collatz Conjecture},
arXiv:2105.14697v3 [cs.LO], 2022.

\bibitem{AntihydraMachine}
Ralf Stephan (with Claude Code), \texttt{AntihydraMachine.lean}: BusyLean
machine-checked macro analysis of the Antihydra BB(6) candidate, 2026.

\bibitem{AntihydraSRS}
Ralf Stephan (with Claude Code), \texttt{AntihydraSRS.lean}: the
mixed-base string-rewriting system for Antihydra and its correctness
proof, 2026.

\bibitem{AntihydraSRSSimulation}
Ralf Stephan (with Claude Code), \texttt{AntihydraSRSSimulation.lean}:
faithful simulation of the Antihydra macro orbit by arbitrary SRS
derivations, 2026.

\bibitem{AntihydraSRSForward}
Ralf Stephan (with Claude Code), \texttt{AntihydraSRSForward.lean}:
constructive forward simulation of the Antihydra macro orbit by an
explicit rewriting strategy, 2026.

\bibitem{AntihydraSRSObstruction}
Ralf Stephan (with Claude Code), \texttt{AntihydraSRSObstruction.lean}:
pumping obstruction to regular invariants of Antihydra over
unary-counter encodings, 2026.

\bibitem{AntihydraSRSaxiom}
Ralf Stephan (with Claude Code), \texttt{AntihydraSRSaxiom.lean}: the
BusyLean-free Antihydra macro model used by \texttt{AntihydraSRS.lean},
2026.

\bibitem{AntihydraMahler}
Ralf Stephan (with Claude Code), \texttt{AntihydraMahler.lean}: the
Mahler-form value coordinate for Antihydra and its equivalence to the
BusyLean-free macro model, 2026.

\bibitem{MixedBase}
Ralf Stephan (with Claude Code), \texttt{MixedBaseRepresentation.lean}:
mixed-base representations and the base-swap lemma, 2026.

\bibitem{BasicSRS}
Ralf Stephan (with Claude Code), \texttt{Basic.lean}: string rewriting
systems, rewrite steps, and termination (Book--Otto \S2), 2026.

\bibitem{ReportB3}
Ralf Stephan (with Claude Code), \textit{Report B3 --- Detailed plans
for every research idea in Deliverable 3}, project internal report,
\S8 ``Formal verification track in Lean 4'', 2026.

\bibitem{ForwardStrategy}
Ralf Stephan (with Claude Code), \textit{3-forward-strategy.html} ---
Forward-strategy analysis of the YAH paper, project internal report,
\S5 ``Formal verification track'', 2026.
\end{thebibliography}

\end{document}
